% Generated by scripts/assemble.py; corrections belong in source/corrections.json.
\examhead{17}{中国科学院大学 2025 年数学分析试卷}{22}
% p023-b01; source PDF page 23
\blockanchor{p023-b01}
\question{1}计算下列极限.\par
% p023-b02; source PDF page 23
\blockanchor{p023-b02}
(1) $\lim_{x\to 1}x^{\frac{1}{1 - x}}$\par
% p023-b03; source PDF page 23
\blockanchor{p023-b03}
(2)  $\lim_{n\to\infty}\frac{1}{n^{4}}(1+2^{3}+\cdots+n^{3}).$\par
% p023-b04; source PDF page 23
\blockanchor{p023-b04}
\question{2}用闭区间套定理证明 [0,1] 不可数.\par
% p023-b05; source PDF page 23
\par\Needspace{6\baselineskip}
\blockanchor{p023-b05}
\question{3}设 $f(x)$ 在 $[0,1]$ 上连续可微, $f(0) = 0$. 证明:\par
\nopagebreak[4]
% p023-b06; source PDF page 23
\blockanchor{p023-b06}
\displayformula{\int_ {0} ^ {1} | f ^ {\prime} (x) f (x) | \mathrm{d} x \leq \frac {1}{2} \int_ {0} ^ {1} | f ^ {\prime} (x) | ^ {2} \mathrm{d} x.}
% p023-b07; source PDF page 23
\par\Needspace{6\baselineskip}
\blockanchor{p023-b07}
\question{4}求重积分\par
\nopagebreak[4]
% p023-b08; source PDF page 23
\blockanchor{p023-b08}
\displayformula{\iint_ {D} (x + y) \sin (x - y) \mathrm{d} x \mathrm{d} y.}
% p023-b09; source PDF page 23
\blockanchor{p023-b09}
其中  $D=\left\{(x,y)\mid0\leq x+y\leq\pi,0\leq x-y\leq\pi\right\}$ .\par
% p023-b10; source PDF page 23
\par\Needspace{6\baselineskip}
\blockanchor{p023-b10}
\question{5}求曲面积分\par
\nopagebreak[4]
% p023-b11; source PDF page 23
\blockanchor{p023-b11}
\displayformula{\iint_ {\Gamma} (y ^ {2} - x) \mathrm{d} y \mathrm{d} z + (z ^ {2} - y) \mathrm{d} z \mathrm{d} x + (x ^ {2} - z) \mathrm{d} x \mathrm{d} y.}
% p023-b12; source PDF page 23
\blockanchor{p023-b12}
其中 $\Gamma = \{(x,y,z) \mid z = 1 - x^2 - y^2, z \geq 0\}$, 方向取外侧.\par
% p023-b13; source PDF page 23
\par\Needspace{6\baselineskip}
\blockanchor{p023-b13}
\question{6}设 $f(x)$ 在 $\left(-\frac{\pi}{2}, 0\right)$ 上可积或绝对可积, 怎么将 $f(x)$ 延拓为到 $(- \pi, \pi)$ 上, 满足其傅里叶级数为\par
\nopagebreak[4]
% p023-b14; source PDF page 23
\blockanchor{p023-b14}
\displayformula{\sum_ {n = 1} ^ {\infty} b _ {2 n - 1} \sin (2 n - 1) x.}
% p023-b15; source PDF page 23
\par\Needspace{6\baselineskip}
\blockanchor{p023-b15}
\question{7}设  $f(x)$  在  $[0, \pi]$  上可积，且  $\int_{0}^{\pi} f(x) \, dx = \pi$ . 求系数  $c_{1}, c_{2}, \cdots, c_{n}$ ，使得\par
\nopagebreak[4]
% p023-b16; source PDF page 23
\blockanchor{p023-b16}
\displayformula{\int_ {0} ^ {\pi} \left[ f (x) - \sum_ {k = 1} ^ {n} c _ {k} \cos k x \right] ^ {2} \mathrm{d} x}
% p023-b17; source PDF page 23
\blockanchor{p023-b17}
最小, 并求 $\lim_{n\to \infty}\sum_{k = 1}^{n}c_k\cos kx = F(x)$ 的表达式.\par
% p023-b18; source PDF page 23
\par\Needspace{6\baselineskip}
\blockanchor{p023-b18}
\question{8}设 $f(x)$ 三次连续可微, 且\par
\nopagebreak[4]
% p023-b19; source PDF page 23
\blockanchor{p023-b19}
\displayformula{f ^ {\prime} (0) = 1, f ^ {\prime \prime} (0) = 0, f ^ {\prime \prime \prime} (0) = - 1.}
% p023-b20; source PDF page 23
\blockanchor{p023-b20}
设 $f(a_{n}) = a_{n + 1}$ ，且 $\lim_{n\to \infty}a_n = 0.$ 求 $\lim_{n\to \infty}n(a_n)^2.$\par
% p023-b21; source PDF page 23
\blockanchor{p023-b21}
\question{9}设 $\{a_{n}\}$ 单调收敛到 $0, \alpha \in (0, \pi)$, 证明: $\sum_{n=1}^{\infty} a_{n} \cos nx$ 在 $[\alpha, 2\pi - \alpha]$ 上一致收敛.\par
% p023-b22; source PDF page 23
\blockanchor{p023-b22}
\question{10}求  $\sum_{n=1}^{\infty}\left(1+\frac{1}{2}+\cdots+\frac{1}{n}\right)x^{n}$  的收敛域与和函数.\par
